% ============ 第2章 §2.20 半平面上的Hardy空间(原书 pp.190-200) ============
\section{半平面上的 Hardy 空间}\label{sec:2.20}
\index{Hardy 空间!(上半平面上的)}
\index{Blaschke 积!(上半平面上的)}
\index{内函数!(上半平面上的)}
\index{外函数!(上半平面上的)}
\index{Poisson 核!(上半平面的)}


$\mathcal{H}(\mathcal{D})$ 表示上半平面 $\mathcal{D} = \{z \in \mathbb{C}; \operatorname{Im}(z) > 0\}$ 上正则函数的全体.对 $f \in \mathcal{H}(\mathcal{D})$,当 $0 < y < \infty$ 时令
\begin{equation*}
\mathcal{M}_{p}(y, f) = \left\{\int_{-\infty}^{\infty}|f(x + iy)|^{p}\dd x\right\}^{1/p} \quad (0 < p < \infty).
\end{equation*}
$\mathcal{D}$ 上的 Hardy 空间 $\mathcal{H}^{p}$ 定义如下:
\begin{equation*}
\mathcal{H}^{p} = \left\{f \in \mathcal{H}(\mathcal{D}); \sup_{0 < y < \infty}\mathcal{M}_{p}(y, f) < \infty \right\},
\end{equation*}
且 $\mathcal{H}^{\infty}$ 表示 $\mathcal{D}$ 上有界正则函数的全体.

本节作为单位圆盘 $D$ 上的 Hardy 空间 $H^{p}$ 的结果的应用,叙述上半平面 $\mathcal{D}$ 上的 Hardy 空间 $\mathcal{H}^{p}$.令 $\overline{\mathcal{D}} = \{z \in \mathbb{C}; \operatorname{Im}(z) \geq 0\}$ 且
\begin{equation*}
z = \phi(w) = \frac{i - w}{i + w} \quad (w \in \overline{\mathcal{D}}),
\end{equation*}
则 $\phi$ 给出从 $\overline{\mathcal{D}}$ 到 $\overline{D}$ 的一对一正则映射.
\begin{equation*}
w = \psi(z) = i\frac{1 - z}{1 + z},
\end{equation*}
则 $\psi \circ \phi(w) = w$ 且 $\phi \circ \psi(z) = z$,$\psi$ 是从 $\overline{D}$ 到 $\overline{\mathcal{D}}$ 的一对一正则映射.对 $g \in H(D)$,$g \circ \phi \in \mathcal{H}(\mathcal{D})$;反之,对 $f \in \mathcal{H}(\mathcal{D})$,$f \circ \psi \in H(D)$.但对 $g \in H^{p}$,$g \circ \phi \in \mathcal{H}^{p}$ 未必成立.实际上,对某个 $0 < \alpha < 1$,$g(z) = 1/(1-z)^{\alpha}$ 属于 $H^{1}$,但容易证明 $g \notin \bigcup_{p \geq 1}\mathcal{H}^{p}$.另一方面,对 $f \in \mathcal{H}^{p}$ 可以证明 $f \circ \psi \in H^{p}$,故 2.19 节为止关于 $H^{p}$ 的结果可用于 $\mathcal{H}^{p}$ 的研究.另外,若 g 在 $D$ 上调和,则 $g \circ \phi$ 在 $\mathcal{D}$ 上调和;反之,若 f 在 $\mathcal{D}$ 上调和,则 $f \circ \psi$ 在 $D$ 上调和.

在单位圆盘 $D$ 上的 Hardy 空间的研究中,Cauchy 核,Poisson 核与 Herglotz 核起了重要作用.把这 3 个核变换到上半平面 $\mathcal{D}$ 上.写 $\phi(t) = e^{i\tau}$($t \in \partial \mathcal{D}$),Cauchy 核变换为
\begin{equation*}
\frac{1}{1 - z} = \frac{1}{2}\left(1 + \frac{i}{w}\right), \quad \frac{e^{i\tau}}{e^{i\tau} - z} = \frac{1}{2}\left(1 + i\frac{1 + tw}{w - t}\right),
\end{equation*}
Poisson 核变换为
\begin{equation*}
\operatorname{Re}\left(\frac{1 + z}{1 - z}\right) = \operatorname{Re}\left(\frac{i}{w}\right), \quad \operatorname{Re}\left(\frac{e^{i\tau} + z}{e^{i\tau} - z}\right) = \operatorname{Re}\left(i\frac{1 + tw}{w - t}\right),
\end{equation*}
Herglotz 核变换为
\begin{equation*}
\frac{1 + z}{1 - z} = \frac{i}{w}, \quad \frac{e^{i\tau} + z}{e^{i\tau} - z} = i\frac{1 + tw}{w - t}.
\end{equation*}
这些变换后的核在 $\mathcal{D}$ 上再次变得重要.

在 $D$ 上重要的一阶变换如下变换.设 $a = \phi(b)$,则
\begin{equation*}
\frac{a - z}{1 - \bar{a}z} = \frac{\phi(b) - \phi(w)}{1 - \overline{\phi(b)}\phi(w)} = \frac{i - \bar{b}}{i + b}\frac{w - b}{w - \bar{b}},
\end{equation*}
且
\begin{equation*}
\frac{|a|}{a}\frac{a - z}{1 - \bar{a}z} = \frac{|\phi(b)|}{\phi(b)}\frac{\phi(b) - \phi(w)}{1 - \overline{\phi(b)}\phi(w)} = \frac{|1 + b^{2}|}{1 + b^{2}}\frac{w - b}{w - \bar{b}}.
\end{equation*}
若设 $\phi(t) = e^{i\tau}$($-\infty < t < \infty$),则对 $g \in L^{1}[-\pi, \pi]$,令 $f(t) = g \circ \phi(t)$,
\begin{equation*}
\frac{1}{2\pi}\int_{-\pi}^{\pi} g(e^{i\tau})\dd \tau = \frac{1}{\pi}\int_{-\infty}^{\infty} f(t)\frac{1}{1 + t^{2}}\dd t.
\end{equation*}
设 $z = re^{i\theta}$,$w = x + iy$ 且 $z = \phi(w)$,$0 \leq r < 1$,则
\begin{equation*}
\frac{1}{2\pi}\int_{-\pi}^{\pi} g(e^{i\tau})P(r, \theta - \tau)\dd \tau = \frac{1}{\pi}\int_{-\infty}^{\infty} f(t)P(y, x - t)\dd t.
\end{equation*}
其中 $P(y, x - t) = y/((x-t)^{2} + y^{2})$.$P(y, x) = \mathrm{Re}(i/w) = y/(x^{2} + y^{2})$ 称为 $\mathcal{D}$ 上的 Poisson 核.

\begin{theorem}\label{thm:2.20.1}
设 F 在 $(-\infty, \infty)$ 上关于 $(1 + t^{2})^{-1}\dd t$ 可积.当 y > 0 时,
\begin{equation*}
f(x + iy) = \frac{1}{\pi}\int_{-\infty}^{\infty} F(t)P(y, x - t)\dd t,
\end{equation*}
则 f 在 $\mathcal{D}$ 上调和.进而 $\lim_{y \to 0}f(x + iy) = F(x)$ a.e. $x \in (-\infty, \infty)$.这里 $y \to 0$ 指的是 $y = (1 - r^{2})/(1 + r^{2} + 2r\cos\theta) \to 0$($r \to 1$).
\end{theorem}

\begin{proof}
令 $G(e^{i\tau}) = (F \circ \psi)(e^{i\tau}) = F(t)$,因 F 关于 $(1 + t^{2})^{-1}\dd t$ 可积,$G \in L^{1}[-\pi, \pi]$.于是设 $z = re^{i\theta}$,$\psi(z) = w = x + iy$,则
\begin{align*}
g(re^{i\theta}) & = \frac{1}{2\pi}\int_{-\pi}^{\pi} G(e^{i\tau})P(r, \theta - \tau)\dd \tau \\
& = \frac{1}{\pi}\int_{-\infty}^{\infty} F(t)P(y, x - t)\dd t = f(x + iy).
\end{align*}
由
\begin{equation*}
y = \frac{1 - r^{2}}{1 + r^{2} + 2r\cos\theta}, \quad x = \frac{2r\sin\theta}{1 + r^{2} + 2r\cos\theta},
\end{equation*}
由定理~\ref{thm:2.3.2} 及其后的注意,$\lim_{r \to 1} g(re^{i\theta}) = G(e^{i\theta})$ a.e. $\theta \in [-\pi, \pi]$,故可证明 $\lim_{y \to 0}f(x + iy) = F(x)$ a.e. $x \in (-\infty, \infty)$.
\end{proof}

已知沿与 x 轴正交的直线取 $y \to 0$ 时定理~\ref{thm:2.20.1} 也成立.

\begin{lemma}\label{lem:2.20.2}
设 $p \geq 1$ 且 $f \in \mathcal{H}^{p}$.设 $w = \zeta + i\eta$ 且 $y_{0} > 0$,则以下成立.
\begin{equation*}
f(w) = \frac{1}{2\pi i}\int_{-\infty}^{\infty}\frac{f(t + iy_{0})}{t + iy_{0} - w}\dd t \quad (\eta > y_{0}), \qquad (1)
\end{equation*}
\begin{equation*}
\frac{1}{2\pi i}\int_{-\infty}^{\infty}\frac{f(t + iy_{0})}{t + iy_{0} - w}\dd t = 0 \quad (\eta < y_{0}), \qquad (2)
\end{equation*}
\begin{equation*}
f(w) = \frac{1}{\pi}\int_{-\infty}^{\infty} f(t + iy_{0})\frac{\eta - y_{0}}{(t - \zeta)^{2} + (\eta - y_{0})^{2}}\dd t \quad (\eta > y_{0}). \qquad (3)
\end{equation*}
\end{lemma}

\begin{proof}
(1) 固定 w,设 $\operatorname{Im}(w) > 0$.选取 x, y 使得 $-x < \mathrm{Re}(w) < x$,$y_{0} < \mathrm{Im}(w) < y$.由对矩形 $-x + iy_{0}$,$x + iy_{0}$,$x + iy$,$-x + iy$ 的 Cauchy 积分公式,
\begin{align*}
2\pi i f(w) & = \int_{-x}^{x}\frac{f(t + iy_{0})}{t + iy_{0} - w}\dd t - \int_{-x}^{x}\frac{f(t + iy)}{t + iy - w}\dd t \\
& \quad + i\int_{y_{0}}^{y}\frac{f(x + is)}{x + is - w}\dd s - i\int_{y_{0}}^{y}\frac{f(-x + is)}{-x + is - w}\dd s \\
& = I_{1}(x) + I_{2}(x) + I_{3}(x) + I_{4}(x),
\end{align*}
取 $\alpha$ 使 $\alpha > 2|\mathrm{Re}(w)|$,将上式两边在区间 $[\alpha, 2\alpha]$ 上平均,得
\begin{equation*}
2\pi i f(w) = \sum_{j=1}^{4}\frac{1}{\alpha}\int_{\alpha}^{2\alpha}I_{j}(x)\dd x.
\end{equation*}
对每个 j 令
\begin{equation*}
I_{j}(\infty) = \lim_{\alpha \to \infty}\frac{1}{\alpha}\int_{\alpha}^{2\alpha}I_{j}(x)\dd x,
\end{equation*}
则 $I_{3}(\infty) = I_{4}(\infty) = 0$(后文证明),
\begin{equation*}
I_{1}(\infty) = \int_{-\infty}^{\infty}\frac{f(t + iy_{0})}{t + iy_{0} - w}\dd t
\end{equation*}
且
\begin{equation*}
I_{2}(\infty) = -\int_{-\infty}^{\infty}\frac{f(t + iy)}{t + iy - w}\dd t
\end{equation*}
可以证明(后文证明).于是对任意 $y > \operatorname{Im}(w)$,
\begin{equation*}
2\pi i f(w) = \int_{-\infty}^{\infty}\frac{f(t + iy_{0})}{t + iy_{0} - w}\dd t - \int_{-\infty}^{\infty}\frac{f(t + iy)}{t + iy - w}\dd t.
\end{equation*}
令 $y \to \infty$,则可证明
\begin{equation*}
-I_{2}(\infty) = \int_{-\infty}^{\infty}\frac{f(t + iy)}{t + iy - w}\dd t \longrightarrow 0 \ \text{(后文证明)},
\end{equation*}
故
\begin{equation*}
2\pi i f(w) = \int_{-\infty}^{\infty}\frac{f(t + iy_{0})}{t + iy_{0} - w}\dd t \quad (\operatorname{Im}(w) > y_{0}),
\end{equation*}
(1) 的证明结束.

下面证明上述假定成立之处.$I_{1}(\infty)$ 与 $I_{2}(\infty)$ 的积分表示,由于 s = $y_{0}$ 或 s = y 时 $y_{0} < \mathrm{Im}(w) < y$,
\begin{equation*}
\int_{-\infty}^{\infty}\frac{|f(t + is)|}{|t + is - w|}\dd t < \infty,
\end{equation*}
由 $\lim_{x \to \infty}I_{j}(x)$(j = 1,2)存在即明.$I_{2}(\infty)$ 的积分表示在 $y \to \infty$ 时收敛到 0,分 p = 1 与 p > 1 两种情形如下证明.p = 1 时,
\begin{align*}
\left| \int_{-\infty}^{\infty}\frac{f(t + iy)}{t + iy - w}\dd t \right| & \leq \int_{-\infty}^{\infty}\frac{|f(t + iy)|}{y - \mathrm{Im}(w)}\dd t \\
& \leq \frac{1}{y - \mathrm{Im}(w)}\sup_{0 < y < \infty}\mathcal{M}_{1}(y, f) = O\left(\frac{1}{y}\right),
\end{align*}
p > 1 时,设 $1/p + 1/q = 1$,由 Hölder 不等式,
\begin{align*}
\left| \int_{-\infty}^{\infty}\frac{f(t + iy)}{t + iy - w}\dd t \right|&\leq \sup_{0 < y < \infty}\mathcal{M}_{p}(y, f)\left(\int_{-\infty}^{\infty}\frac{\dd t}{|t + iy - w|^{q}}\right)^{1/q} = O(y^{-1 + \frac{1}{q}}).
 \end{align*}
至此,还剩下证明假定的 $I_{3}(\infty) = I_{4}(\infty) = 0$.$I_{4}(\infty) = 0$ 类似,只证明 $I_{3}(\infty) = 0$.因 $\alpha > 2|\mathrm{Re}(w)|$,$\alpha < x < 2\alpha$,故 $|x + is - w| \geq |x - \mathrm{Re}(w)| \geq \alpha/2$,所以
\begin{align*}
\left| \frac{1}{\alpha}\int_{\alpha}^{2\alpha}I_{3}(x)\dd x \right| & \leq \frac{1}{\alpha}\int_{\alpha}^{2\alpha}\dd x\int_{y_{0}}^{y}\left| \frac{f(x + is)}{x + is - w} \right|\dd s \\
& \leq \frac{2}{\alpha^{2}}\int_{y_{0}}^{y}\dd s\int_{\alpha}^{2\alpha}|f(x + is)|\dd x.
\end{align*}
设 $1/p + 1/q = 1$,用 Hölder 不等式,
\begin{equation*}
\int_{\alpha}^{2\alpha}|f(x + is)|\dd x \leq \alpha^{1/q}\sup_{0 < y < \infty}\mathcal{M}_{p}(y, f),
\end{equation*}
代入上式得
\begin{equation*}
\left| \frac{1}{\alpha}\int_{\alpha}^{2\alpha}I_{3}(x)\dd x \right| \leq 2\alpha^{\frac{1}{q} - 2}(y - y_{0})\sup_{0 < y < \infty}\mathcal{M}_{p}(y, f).
\end{equation*}
因 $f \in \mathcal{H}^{p}$,令 $\alpha \to \infty$,
\begin{equation*}
I_{3}(\infty) = \lim_{\alpha \to \infty}\frac{1}{\alpha}\int_{\alpha}^{2\alpha}I_{3}(x)\dd x = 0.
\end{equation*}

(2) 因 $\operatorname{Im}(w) < y_{0}$,对矩形 $-x + iy_{0}$,$x + iy_{0}$,$x + iy$,$-x + iy$ 用 Cauchy 定理,
\begin{align*}
0 &= \int_{-x}^{x}\frac{f(t + iy_{0})}{t + iy_{0} - w}\dd t - \int_{-x}^{x}\frac{f(t + iy)}{t + iy - w}\dd t \\
&\quad + i\int_{y_{0}}^{y}\frac{f(x + is)}{x + is - w}\dd s - i\int_{y_{0}}^{y}\frac{f(-x + is)}{-x + is - w}\dd s \\
\quad &= I_{1}(x) + I_{2}(x) + I_{3}(x) + I_{4}(x).
\end{align*}
与 (1) 相同,可证明 $\lim_{y \to \infty}I_{2}(\infty) = I_{3}(\infty) = I_{4}(\infty) = 0$,故
\begin{equation*}
\frac{1}{2\pi i}\int_{-\infty}^{\infty}\frac{f(t + iy_{0})}{t + iy_{0} - w}\dd t = 0 \quad (\operatorname{Im}(w) < y_{0}).
\end{equation*}

(3) 设 $w = \zeta + i\eta$,$\eta > y_{0}$ 且 $w' = i2y_{0} + \zeta - i\eta$.因 $\operatorname{Im}(w') < y_{0}$,由 (2),
\begin{equation*}
0 = \frac{1}{2\pi i}\int_{-\infty}^{\infty}\frac{f(t + iy_{0})}{t + iy_{0} - w'}\dd t,
\end{equation*}
从 (1) 中 $f(w)$ 的积分表示减去它,
\begin{equation*}
f(w) = \frac{1}{\pi}\int_{-\infty}^{\infty} f(t + iy_{0})\frac{\eta - y_{0}}{(t - \zeta)^{2} + (\eta - y_{0})^{2}}\dd t.
\end{equation*}
\end{proof}

\begin{lemma}\label{lem:2.20.3}
设 $1 \leq p < \infty$.设 $\Gamma = \{z; |z| = R < 1\}$ 且 $\gamma$ 为 $\Gamma$ 内部任意圆周,若 $g \in H(D)$,则
\begin{equation*}
\int_{\gamma}|g(z)|^{p}|\dd z| \leq 2\int_{\Gamma}|g(z)|^{p}|\dd z|.
\end{equation*}
\end{lemma}

\begin{proof}
当 $z \in \gamma$ 时,
\begin{equation*}
g(z) = \frac{1}{2\pi}\int_{-\pi}^{\pi} g(Re^{i\theta})\frac{R^{2} - |z|^{2}}{|Re^{i\theta} - z|^{2}}\dd \theta,
\end{equation*}
用 Hölder 不等式,
\begin{align*}
\int_{\gamma}|g(z)|^{p}|\dd z|&\leq \int_{\gamma}\left(\frac{1}{2\pi}\int_{-\pi}^{\pi}|g(Re^{i\theta})|^{p}\frac{R^{2} - |z|^{2}}{|Re^{i\theta} - z|^{2}}\dd \theta\right)|\dd z|  \\
&\leq \int_{-\pi}^{\pi}|g(Re^{i\theta})|^{p}\left(\frac{1}{2\pi}\int_{\gamma}\frac{R^{2} - |z|^{2}}{|Re^{i\theta} - z|^{2}}|\dd z|\right)\dd \theta.
\end{align*}
设圆 $\gamma$ 的中心为 a,半径为 r,则 $r \leq |Re^{-i\theta} - a|$,由 Gauss 平均值定理,
\begin{equation*}
\frac{1}{2\pi}\int_{\gamma}\frac{R^{2} - |z|^{2}}{|Re^{i\theta} - z|^{2}}|\dd z| = r \times \frac{R^{2} - |a|^{2}}{|Re^{i\theta} - a|^{2}} \leq 2.
\end{equation*}
引理得证.
\end{proof}

\begin{theorem}\label{thm:2.20.4}
设 $1 \leq p \leq \infty$.对 $f \in \mathcal{H}^{p}$ 令 $g = f \circ \psi$,则 $g \in H^{p}$.其中 $\psi(z) = i(1 - z)/(1 + z)$.
\end{theorem}

\begin{proof}
p = $\infty$ 时显然,设 $1 \leq p < \infty$.设 $g = f \circ \psi$,若能证明
\begin{equation*}
\sup_{r}\int_{-\pi}^{\pi}|g(re^{i\theta})|^{p}\dd \theta/2\pi \leq 2\sup_{y}\int_{-\infty}^{\infty}|f(x + iy)|^{p}\dd x/\pi,
\end{equation*}
则由假设 $f \in \mathcal{H}^{p}$,g 属于 $H^{p}$.

对 0 < r < 1 令 $\gamma_{r} = \{z; |z| = r\}$,对 $0 < \delta$ 令
\begin{equation*}
\Gamma_{\delta} = \left\{z; z = \phi(w) = \frac{i - w}{i + w}, \operatorname{Im}(w) = \delta \right\},
\end{equation*}
则 $\Gamma_{\delta}$ 是中心 $-\delta/(1+\delta)$,半径 $1/(1+\delta)$ 的圆.对任意 r < 1,存在 $\delta > 0$,使得 $\gamma_{r}$ 包含在 $\Gamma_{\delta}$ 的内部.由引理~\ref{lem:2.20.2} 的 (3),f 在 $\{w; \operatorname{Im}(w) \geq \delta\} \cup \{\text{无穷远点}\}$ 上连续.因为 $w \to \infty$ 时,在 $\{w; \operatorname{Im}(w) \geq \delta\}$ 上 $f(w) \to 0$.于是 g 在圆盘 $\{z; z = \phi(w), \operatorname{Im}(w) \geq \delta\}$ 上连续且在其内部正则.因 $\gamma_{r}$ 包含于该圆盘内部,由引理~\ref{lem:2.20.3} 定理得证.
\end{proof}

\begin{theorem}\label{thm:2.20.5}
设 $1 \leq p < \infty$ 且 $f \in \mathcal{H}^{p}$.$y \to 0$ 的含义与定理~\ref{thm:2.20.1} 相同.

(1) $\lim_{y \to 0}f(x + iy)$ 在 a.e. $x \in (-\infty, \infty)$ 存在.

(2) $f(x) = \lim_{y \to 0}f(x + iy)$ 属于 $L^{p}(-\infty, \infty)$,且
\begin{equation*}
f(x + iy) = \frac{1}{\pi}\int_{-\infty}^{\infty} f(t)P(y, x - t)\dd t \quad (y > 0).
\end{equation*}
\end{theorem}

\begin{proof}
(1) 令 $g = f \circ \psi$,由定理~\ref{thm:2.20.4},$g \in H^{p}$,由定理~\ref{thm:2.5.3},$\lim_{r \to 1}g(re^{i\theta})$ 在 a.e. $\theta \in [-\pi, \pi]$ 存在.于是定理~\ref{thm:2.20.1} 的证明证明了 (1).

(2) 再令 $g = f \circ \psi$,因 $p \geq 1$,由定理~\ref{thm:2.9.1},
\begin{equation*}
g(z) = \frac{1}{2\pi}\int_{-\pi}^{\pi} g(e^{i\tau})P(r, \theta - \tau)\dd \tau .
\end{equation*}
由定理~\ref{thm:2.20.1} 前的注意,设 $\psi(z) = x + iy$,则
\begin{equation*}
f(x + iy) = \frac{1}{\pi}\int_{-\infty}^{\infty} f(t)P(y, x - t)\dd t.
\end{equation*}
\end{proof}

利用定理~\ref{thm:2.20.4},与定理~\ref{thm:2.20.5} 一样,$p \geq 1$ 时在 $H^{p}$ 中成立的许多结果在 $\mathcal{H}^{p}$ 中也成立.这里叙述属于 $\mathcal{H}^{p}$ 的函数的零点.

设
\begin{equation*}
\{b_{j}\}_{j=1}^{\infty} \subset \mathcal{D} \quad \text{且} \quad \sum_{j=1}^{\infty}\frac{\mathrm{Im}(b_{j})}{|i + b_{j}|^{2}} < \infty,
\end{equation*}
则
\begin{equation*}
\mathcal{B}(w) = \left(\frac{w - i}{w + i}\right)^{m}\prod_{j=1}^{\infty}\frac{|b_{j}^{2} + 1|}{b_{j}^{2} + 1}\cdot\frac{w - b_{j}}{w - \bar{b}_{j}}
\end{equation*}
称为上半平面 $\mathcal{D}$ 上的 Blaschke 积.这里 m 是非负整数,且允许 $b_{j} = b_{i}$($j \neq i$).设 $a_{j} = \phi(b_{j})$,$z = \phi(w)$,则
\begin{equation*}
\frac{|a_{j}|}{a_{j}}\frac{a_{j} - z}{1 - \bar{a}_{j}z} = \frac{|b_{j}^{2} + 1|}{b_{j}^{2} + 1}\cdot\frac{w - b_{j}}{w - \bar{b}_{j}},
\end{equation*}
故令 $B(z) = \mathcal{B} \circ \psi(z)$,即 $B \circ \phi(w) = \mathcal{B}(w)$,则 B(z) 是 D 上 $\{a_{j}\}_{j=1}^{\infty}$ 的 Blaschke 积.实际上
\begin{equation*}
1 - |a_{j}|^{2} = 1 - \left| \frac{i - b_{j}}{i + b_{j}} \right|^{2} = \frac{4\mathrm{Im}(b_{j})}{|i + b_{j}|^{2}},
\end{equation*}
且由
\begin{equation*}
\sum_{j=1}^{\infty}\frac{\operatorname{Im}(b_{j})}{|i + b_{j}|^{2}} < \infty \quad \text{得} \quad \sum_{j=1}^{\infty}(1 - |a_{j}|^{2}) < \infty.
\end{equation*}
由定理~\ref{thm:2.6.1},$B \in H^{\infty}$,故 $\mathcal{B} \in \mathcal{H}^{\infty}$,$|\mathcal{B}(w)| < 1$($w \in \mathcal{D}$)且 $|\mathcal{B}(x)| = 1$ a.e. $x \in (-\infty, \infty)$.$\{b_{j}\}_{j=1}^{\infty}$ 是 $\mathcal{B}$ 的零点,且 $\mathcal{B}$ 在 $\mathcal{D}$ 上没有其他零点.

\begin{theorem}\label{thm:2.20.6}
设 $1 \leq p \leq \infty$.设 $\{b_{j}\}_{j=1}^{\infty}$ 为 $f \in \mathcal{H}^{p}$ 的零点集合,则
\begin{equation*}
\sum_{j=1}^{\infty}\frac{\operatorname{Im}(b_{j})}{|i + b_{j}|^{2}} < \infty.
\end{equation*}
反之,若 $\sum_{j=1}^{\infty}\mathrm{Im}(b_{j})/|i + b_{j}|^{2} < \infty$,则存在以 $\{b_{j}\}_{j=1}^{\infty}$ 为零点集合的 $f \in \mathcal{H}^{p}$.
\end{theorem}

\begin{proof}
令 $g = f \circ \psi$,由定理~\ref{thm:2.20.4},$g \in H^{p}$,设 $a_{j} = \phi(b_{j})$,则 $\{a_{j}\}_{j=1}^{\infty}$ 是 g 的零点集合.由定理~\ref{thm:2.6.3},$\sum_{j=1}^{\infty}(1 - |a_{j}|) < \infty$,故 $\sum_{j=1}^{\infty}\mathrm{Im}(b_{j})/|i + b_{j}|^{2} < \infty$.反之,由本定理上面的注意,设 $\mathcal{B}$ 为 $\{b_{j}\}_{j=1}^{\infty}$ 的 Blaschke 积,则 $\mathcal{B} \in \mathcal{H}^{\infty}$.令 $h(w) = (w + i)^{-2/p}$,则 h 属于 $\mathcal{H}^{p}$ 且在 $\mathcal{D}$ 上没有零点.f = $\mathcal{B}h$ 即为所求.
\end{proof}

\begin{theorem}\label{thm:2.20.7}
设 $1 \leq p \leq \infty$.若 $f \in \mathcal{H}^{p}$ 且 $f(w) \neq 0$,则 $f(w) = \mathcal{B}(w)F(w)$ 可分解.其中 $\mathcal{B}$ 是 Blaschke 积,F 属于 $\mathcal{H}^{p}$,且在 $\mathcal{D}$ 上没有零点.
\end{theorem}

\begin{proof}
因 $g = f \circ \psi \in H^{p}$,由定理~\ref{thm:2.6.4},g = Bh 可分解,其中 B 是 Blaschke 积,$h \in H^{p}$ 在 D 上没有零点.于是由定理~\ref{thm:2.9.1},设 $z = re^{i\theta}$,则
\begin{equation*}
h^{p}(z) = \frac{1}{2\pi}\int_{-\pi}^{\pi} h^{p}(e^{i\tau})P(r, \theta - \tau)\dd \tau
\end{equation*}
可作 Poisson 积分表示.$B \circ \phi$ 是 $\mathcal{D}$ 上的 Blaschke 积,$F = h \circ \phi$ 在 $\mathcal{D}$ 上没有零点.剩下的问题是证明 $F \in \mathcal{H}^{p}$.由定理~\ref{thm:2.20.1} 前的注意,
\begin{equation*}
F^{p}(x + iy) = \frac{1}{\pi}\int_{-\infty}^{\infty} F^{p}(t)P(y, x - t)\dd t \quad (y > 0),
\end{equation*}
利用
\begin{equation*}
\frac{1}{\pi}\int_{-\infty}^{\infty} P(y, x - t)\dd x = \frac{1}{2\pi}\int_{-\pi}^{\pi} P(r, \theta - \tau)\dd \tau = 1,
\end{equation*}
由 Fubini 定理,
\begin{equation*}
\sup_{y > 0}\int_{-\infty}^{\infty}|F(x + iy)|^{p}\dd x < \infty,
\end{equation*}
故可证明 $F \in \mathcal{H}^{p}$.这里用到了 $|F(t)| = |f(t)|$ a.e. t.
\end{proof}

对关于 $(-\infty, \infty)$ 上 Lebesgue 测度奇异的(有界)Borel 正测度 $\nu$,
\begin{equation*}
\mathcal{S}(w) = \exp\left\{i\int_{-\infty}^{\infty}\frac{1 + tw}{t - w}\dd \nu(t)\right\}
\end{equation*}
称为奇异内函数.已知 $\mathcal{D}$ 上的任意内函数(即属于 $\mathcal{H}^{\infty}$,其边界值几乎处处绝对值等于 1 的函数)存在常数 $\alpha \geq 0$,可写成
\begin{equation*}
e^{i\alpha w}\mathcal{B}(w)\mathcal{S}(w).
\end{equation*}
若存在实常数 $\gamma$,$G \in \mathcal{H}^{p}$ 可写成
\begin{equation*}
G(w) = e^{i\gamma}\exp\left\{\frac{1}{\pi i}\int_{-\infty}^{\infty}\frac{1 + tw}{t - w}\frac{\log|G(t)|}{1 + t^{2}}\dd t \right\},
\end{equation*}
则称为 $\mathcal{D}$ 上的外函数.已知任意 $\mathcal{H}^{p}$ 的函数可写成内函数与外函数的乘积.
